\documentclass[11pt]{article}
\usepackage[margin=1in]{geometry}
\usepackage{amsmath,amssymb,amsthm,hyperref}
\newtheorem{theorem}{Theorem}
\title{A colored-comb basis for the full weighted partition poset}
\author{Wentao Li}
\date{November 20, 2026}
\begin{document}\maketitle
\begin{abstract}
Gonz\'alez D'Le\'on conjectured that colored combs whose root avoids the largest
color index a basis for the top reduced cohomology of the full weighted
partition poset with its minimum removed. He proved they span. Decomposing
a comb at a root of the largest color gives a differential equation for the
remaining combs. Its solution agrees with the Euler characteristic of the
full poset, proving independence for every number of colors.
\end{abstract}

Fix $k\ge1$. A normalized binary tree on distinct ordered labels is a full
plane binary tree in which the smallest leaf of each subtree is its leftmost
leaf. Color each internal vertex in $[k]$. A colored comb requires that
whenever the right child is internal its color be strictly smaller than its
parent's. Let $C_n$ count all such combs on $[n]$ and let $D_n$ count those
whose root color differs from $k$. Count the one-leaf tree in both classes.

In $\Pi_n^k$ each block carries a weak $k$-composition of its size minus one.
Merging two blocks adds their weights and a unit vector of the chosen color.
The chain of a comb successively merges the two child blocks at each internal
node, in postorder and with its specified color. Delete the bottom element
from this chain to get the cohomology class denoted by $\check c$ in
\cite{g}. Cohomology is over the ground field of the source.
\begin{theorem}
For $n\ge2$, the classes $\check c$ of the colored combs whose root color
differs from $k$ form a basis of
$\widetilde H^{n-2}(\Pi_n^k\setminus\{\hat0\})$.
\end{theorem}
\begin{proof}
Set $C(z)=\sum_{n\ge1} C_nz^n/n!$ and
$D(z)=\sum_{n\ge1}D_nz^n/n!$. A comb counted by $C_n-D_n$ has root color
$k$. Its left subtree is any comb containing the smallest label, and its
right subtree is precisely a comb counted by $D$, including a single leaf.
Conversely these data determine a unique such comb. Thus
\[
 C_n-D_n=\sum_{a=1}^{n-1}\binom{n-1}{a-1} C_aD_{n-a}.
\]
It follows that $C'-D'=C'D$. Both series have zero constant term, so
\[
                         D(z)=1-e^{-C(z)}.                       \tag{1}
\]

Write $\widehat\Pi_n^k$ for the poset after adjoining a new top.
For a weighted block of size $s$ sum the bottom-to-block M\"obius values over
all its weights and denote the answer by $b_s$. The sign and comb-basis
theorems for the \emph{one-block intervals} in \cite{g} (Theorem 1.3 and
Corollary 5.6, or Theorem 5.7) give $b_s=(-1)^{s-1}C_s$ for $s\ge2$.
For $s=1$ both sides are one. Bottom intervals factor independently over
the blocks. Therefore if $M_n=\sum_{\pi\in\Pi_n^k}\mu(\hat0,\pi)$, the
labelled partition formula gives
\[
       1+\sum_{n\ge1}M_nz^n/n!
          =\exp\left(\sum_{s\ge1}b_s z^s/s!\right)
          =\exp(-C(-z)).                                         \tag{2}
\]
The new top has M\"obius value $-M_n$. The proper part of the augmented
poset is exactly the poset in the theorem. By EL-shellability
(\cite{g}, Theorem 1.4) its reduced cohomology is concentrated in degree
$n-2$. Its dimension $d_n$ is consequently $(-1)^{n-1}M_n$ by the
M\"obius--Euler characteristic identity. Take coefficients in (1) and (2)
to obtain $D_n=d_n$. Proposition 5.9 of \cite{g} proves these $D_n$ classes
span. A spanning family of that size is a basis.
\end{proof}
For $n=2$ the poset above the minimum has $k$ points. Exactly $k-1$ combs
survive the restriction, as expected. For $k=1$ and $n\ge2$ the proposed
family is empty and the proper part under consideration has a maximum.
Both conventions agree with the proof.

\paragraph{Status and credit.} This proves Conjecture 5.11 in the 2014
preprint \cite{g}, also posed in the author's dissertation as Conjecture 6.3.4.
The one-block comb basis, the spanning theorem for the full poset, and the
two-color case are prior results. The new step is the count after excluding
the last root color. Research draft prepared with assistance from OpenAI
Codex. Independent proof and priority review is pending.
\begin{thebibliography}{9}
\bibitem{g} R. S. Gonz\'alez D'Le\'on, \emph{On the free Lie algebra with
multiple brackets}, Adv. Appl. Math. 79 (2016), 37--97. Numbering above
follows the 2014 preprint. \url{https://arxiv.org/abs/1408.5415}
\bibitem{d} R. S. Gonz\'alez D'Le\'on, \emph{On the Combinatorics of the
Free Lie Algebra with Multiple Brackets}, University of Miami, 2014,
Section 6.3. \url{https://scholarship.miami.edu/esploro/outputs/doctoral/On-the-Combinatorics-of-the-Free/991031448073702976}
\end{thebibliography}
\end{document}
